\section{Условие задачи}

\subsection{Цель работы}

Реализовать класс множества объектов с ручным управлением динамической памятью, применить правило пяти и перегрузку операторов. Проверить работу контейнера модульными тестами Google Test.

\subsection{Вариант 19}

Элемент \texttt{ExamGrade} содержит id студента и оценку. Множество \texttt{ExamGradeSet} хранит не более одной оценки на студента. Метод \texttt{add} при совпадении id заменяет оценку и возвращает \texttt{true}, если состояние изменилось. Объединение и пересечение выполняются по id.

Необходимо реализовать:
\begin{itemize}
  \item конструктор по умолчанию, копирование, перемещение и деструктор;
  \item копирующее присваивание через copy-and-swap и перемещающее присваивание;
  \item методы \texttt{add}, \texttt{remove}, \texttt{contains}, \texttt{size}, \texttt{empty};
  \item операторы принадлежности \texttt{()}, добавления \texttt{+=}, удаления \texttt{-=};
  \item объединение, пересечение и разность множеств;
  \item сравнение множеств без учёта порядка хранения и вывод в поток.
\end{itemize}

\subsection{Требования к реализации}

Хранение элементов осуществляется в динамическом массиве \texttt{ExamGrade*} с полями размера и вместимости. Управление памятью выполняется через \texttt{new[]} и \texttt{delete[]}. Запрещены шаблоны, контейнеры STL, умные указатели, \texttt{std::string}, \texttt{std::array}, \texttt{std::span} и \texttt{std::optional}. Используются C++20, CMake 3.20+ и Google Test.
